%!TeX root=csp-proof.tex
\section{The main statements}\label{sec:main-statements}

Everything converges here.
Theorem~\ref{thm:main-inductive} is the single induction that carries the whole
proof; the three theorems the algorithm consumes are corollaries of it.

\subsection{The three claims the algorithm needs}

The algorithm of \S\ref{sec:algorithm} is justified by exactly three facts.
Informally:

\begin{description}\tightlist
\item[\textup{(IC1)}]\label{ic:existence} every domain of size $\ge 2$ has a
  strong subset, or an equivalence $\sigma$ with
  $D_{x}/\sigma\cong\mathbb Z_{p}$;
\item[\textup{(IC2)}]\label{ic:safe} reducing a variable's domain to a strong
  subset of it does not destroy solvability, for a consistent enough instance;
\item[\textup{(IC3)}]\label{ic:codim} for a linked, consistent enough,
  strong-subset-free instance, the set of ``linear cosets'' containing a
  solution is empty, full, or a codimension-one affine subspace.
\end{description}

(IC1) is Lemma~\ref{lem:ubiquity} together with
Lemma~\ref{lem:linear-on-top}. (IC2) is
Theorem~\ref{thm:reductions-safe}. (IC3) is
Theorem~\ref{thm:codim-one}.

\subsection{The main induction}

\begin{longtheorem}[Main inductive claim]\label{thm:main-inductive}
Let $\mathcal I$ be an irreducible cycle-consistent instance and let $D^{(1)}$
be a $1$-consistent reduction for $\mathcal I$ with $D^{(1)}\lll D$.
\begin{enumerate}\tightlist
\item[\textup{(1)}] If $\mathcal I$ is crucial in $D^{(1)}$ then
  \textup{(1a)} holds, and \textup{(1b)} or \textup{(1c)} holds:
  \begin{enumerate}\tightlist
  \item[\textup{(1a)}] every constraint of $\mathcal I$ has the parallelogram
    property;
  \item[\textup{(1b)}] $\mathcal I$ is a connected linear-type instance whose
    solution set is subdirect;
  \item[\textup{(1c)}] there is an expanded covering $\mathcal J$ of
    $\mathcal I$, crucial in $D^{(1)}$, with a linked connected subinstance
    $\Upsilon$ whose solution set is not subdirect.
  \end{enumerate}
\item[\textup{(2)}] If $D^{(2)}\subte{T}D^{(1)}$ is a $1$-consistent reduction
  for $\mathcal I$ with $T\in\{\TBA,\TC\}$ --- and, when $T=\TBA$, uniform:
  $D^{(2)}\subte{\TBA(t)}D^{(1)}$ for a single binary term $t$
  (Caveat~\ref{haz:uniform-witness}) --- and $\mathcal I^{(1)}$ has a solution,
  then $\mathcal I^{(2)}$ has a solution.
\end{enumerate}
\end{longtheorem}

\begin{hazard}[The induction, stated precisely]\label{haz:the-induction}
``Induction on the size of $D^{(1)}$'' hides the entire well-foundedness
argument, and it must be made precise, because \emph{the two halves of the
conclusion are proved by mutual induction and each invokes the other}.

The intended measure is
\[
  \mu(D^{(1)})=\sum_{x\in\Var(\mathcal I)}\bigl|D^{(1)}_{x}\bigr|\in\mathbb N,
\]
and the statement being proved by strong induction on $\mu$ is the
conjunction of (1) and (2), \emph{universally quantified over $\mathcal I$}:
\[
  P(k):\quad
  \forall\,\mathcal I,\ \forall\,D^{(1)}\ \text{with}\ \mu(D^{(1)})\le k,\
  \textup{(1)}\ \text{and}\ \textup{(2)}\ \text{hold}.
\]
The quantifier over $\mathcal I$ must be \emph{inside}, because every use of the
inductive hypothesis changes the instance: (2) applies it to a weakening
$\mathcal I'$ of $\mathcal I^{(2)}$; Case~1 of (1) applies it to $\mathcal I$
with $D^{(2)}$; Case~2 applies it to weakenings and to expanded coverings. An
induction that fixed $\mathcal I$ and inducted on $D^{(1)}$ alone does not go
through.

Three further points.
\begin{enumerate}\alphlabels\tightlist
\item \emph{The two halves are established in a fixed order at each level.}
  Part (2) at a given measure uses (1) only at strictly smaller measure,
  whereas part (1) at that measure uses \emph{(2) at the same measure} ---
  Case~1 below does exactly that. So the scheme is: strong induction on
  $k=\mu(D^{(1)})$; at each level establish (2) for all instances, then (1) for
  all instances. That is acyclic, and it is why part (2) is proved first.
\item \emph{Which appeals decrease.} (2) invokes the hypothesis at
  $D^{(2)}\subsetneq D^{(1)}$. The containment is strict in the only case that
  needs an argument: $D^{(2)}\subte{T}D^{(1)}$ permits equality, but if
  $D^{(2)}=D^{(1)}$ then $\mathcal I^{(2)}=\mathcal I^{(1)}$ and (2) is its own
  hypothesis. Case~1 of (1) likewise. Case~2 invokes it at
  $D^{(B,\bot)}\subseteq D^{(B,\top)}$, which is $D^{(1)}$ with the $z$-domain
  cut to $B\lneq D^{(1)}_{z}$, the properness being that of
  $B\subt[D_{z}]{T(\sigma)}D^{(1)}_{z}$. In every case the decrease is at a
  single variable and by at least one element. The uniform statement is:
  \emph{if $D'\subte{T}D$ with $T$ any type and $D'\neq D$, then
  $\mu(D')<\mu(D)$}.
\item \emph{The measure does survive, because no appeal is made at an
  expanded covering.} This looks like the weak point of the scheme and is not.
  An expanded covering has \emph{more} variables than $\mathcal I$, and the
  reduction of a covering is the extension of $D^{(1)}$ along the parent map
  (Proposition~\ref{prop:covering-basics}(h)), which repeats the parent
  domains; its $\mu$ is \emph{larger}, and no observation about one parent
  variable being cut compares the two sums, since arbitrarily many children can
  outweigh the decrease. So if the induction did appeal at a covering, $\mu$
  would not serve.

  It does not. Every appeal above is to a \emph{weakening} of the current
  instance --- $\Var$ can only shrink --- paired with a reduction strictly
  smaller in $\mu$: part (2) at $D^{(2)}$; Case~1 of (1) at $D^{(2)}$ twice;
  the proof of (1a) at the minimal $\MT T$ reduction; and the three appeals in
  the $\mathcal B_{0}\neq\varnothing$ branch at $D^{(B,\bot)}$, whose members
  $\mathcal I'$ come from $\Omega$, every member of which is a weakening of
  $\mathcal I$ by condition~1 of Construction~\ref{constr:omega}. Expanded
  coverings enter only as \emph{outputs} --- the $\mathcal J$ of (1c), and the
  auxiliary $\bot(\cdot)$, $\Delta(\cdot)$ and $\Theta$ built from them --- and
  an output carries no well-foundedness obligation.

  What does apply this theorem to a covering is
  Theorem~\ref{thm:pc-safe}, which is a separate statement proved
  \emph{after} it and consuming it as a black box; the appeal there is an
  ordinary application of a proved theorem, not a step in this induction.
  Confusing the two is the natural mistake, and it is what an earlier draft of
  this caveat made.
\end{enumerate}
\end{hazard}

\begin{proof}[Proof of \textup{(2)}]
Suppose $\mathcal I^{(2)}$ has no solution. Weaken $\mathcal I^{(2)}$ to an
instance $\mathcal I'$ crucial in $D^{(2)}$ (Remark~\ref{rem:get-crucial}). By
the inductive hypothesis applied to $\mathcal I'$ and $D^{(2)}$, $\mathcal I'$
satisfies (1a) and (1b) or (1c).

\emph{Case \textup{(1c)}.} There is an expanded covering $\mathcal J$ of
$\mathcal I'$, crucial in $D^{(2)}$, with a linked connected subinstance
$\Upsilon$. Let $x$ be a variable of a constraint $C\in\Upsilon$. By
Lemma~\ref{lem:connected}(p), $\ConOne(C,x)$ is a perfect linear congruence;
choose $\zeta\subseteq\mathbf D_{x}\times\mathbf D_{x}\times\mathbf Z_{p}$ with
$(y_{1},y_{2},0)\in\zeta\iff(y_{1},y_{2})\in\ConOne(C,x)$ and
$\proj_{1,2}(\zeta)=\cover{\ConOne(C,x)}$.

Replace the variable $x$ of $C$ in $\mathcal J$ by a fresh $x'$ and add the
constraint $\zeta(x,x',z)$ with $z$ fresh of domain $\mathbf Z_{p}$; call the
result $\Theta$, and extend $D^{(1)},D^{(2)}$ by $D^{(1)}_{x'}=D^{(2)}_{x'}
=D_{x'}$. Let $E_{i}$ be the set of $b\in\mathbf Z_{p}$ such that $\Theta$ has
a solution in $D^{(i)}$ with $z=b$. By
Lemma~\ref{lem:ba-center-implies}, $E_{2}\dsubte{T}E_{1}$.

Since $\mathcal J$ is crucial in $D^{(2)}$, $\Theta^{(2)}$ has a solution but
none with $z=0$; since $\mathcal I^{(1)}$ has a solution, so does
$\mathcal J^{(1)}$, hence $\Theta^{(1)}$ has one with $z=0$. Thus
$0\in E_{1}$, $0\notin E_{2}$, and $E_{1}$ has at least two elements. As
$\mathbf Z_{p}$ has no proper subalgebra of size $>1$, $E_{1}=\mathbf Z_{p}$;
then $E_{2}\subt{T}\mathbf Z_{p}$ with $T\in\{\TBA,\TC\}$, contradicting
Lemma~\ref{lem:no-abs-in-linear}.

\emph{Case \textup{(1b)}.} Corollary~\ref{cor:same-type} gives a
contradiction directly.
\end{proof}

\begin{hazard}[Three gaps in the case \textup{(1c)} argument]\label{haz:case1c}
\begin{enumerate}\alphlabels\tightlist
\item ``$E_{1}$ contains at least two different elements, $0\in E_{1}$ and
  $0\notin E_{2}$, hence $E_{1}=\mathbf Z_{p}$'' uses that $E_{1}$ is a
  subuniverse of $\mathbf Z_{p}$ --- true because $E_{1}=\Theta^{(1)}(z)$ is a
  relation defined by an instance (Definition~\ref{def:instance-relation}) ---
  \emph{and} that the subuniverses of $\mathbf Z_{p}$ are the singletons and
  the whole set. The latter needs $w^{\mathbf Z_{p}}$ to generate: from $a\ne b$
  one reaches every element. Prove it; it is one line but it is the crux.
\item $E_{2}\dsubte{T}E_{1}$ is applied with the \emph{dotted} conclusion of
  Lemma~\ref{lem:ba-center-implies}, and the argument then treats $E_{2}$ as a
  proper subuniverse. The case $E_{2}=\varnothing$ must be handled: it is
  excluded because $\Theta^{(2)}$ has a solution. Say so.
\item The choice of $\zeta$ is by Definition~\ref{def:perfect-linear}, which
  supplies $\zeta$ with $\proj_{1,2}\zeta=\cover\sigma$ and
  $(a_{1},a_{2},b)\in\zeta\Rightarrow((a_{1},a_{2})\in\sigma\iff b=0)$. The
  proof uses the \emph{biconditional} $(y_{1},y_{2},0)\in\zeta\iff(y_{1},y_{2})
  \in\ConOne(C,x)$, which is stronger: it needs, for each $(y_1,y_2)\in\sigma$,
  that $(y_{1},y_{2},0)\in\zeta$. This follows from $\proj_{1,2}\zeta
  =\cover\sigma\supseteq\sigma$ plus the implication, but the step is
  suppressed.
\end{enumerate}
\end{hazard}

\begin{proof}[Proof of \textup{(1)}, sketch of the structure]
First, $|D^{(1)}_{x}|>1$ for some $x$: otherwise $\mathcal I^{(1)}$ has all
domains singletons and $1$-consistency makes it solvable, contradicting
cruciality.

\emph{Case 1: some $D^{(1)}_{x}$ has a nontrivial $\TBA$ or central
subuniverse.} Lemma~\ref{lem:find-consistent} yields a $1$-consistent reduction
$D^{(2)}\subte{T}D^{(1)}$, $T\in\{\TBA,\TC\}$. One checks $\mathcal I$ is
crucial in $D^{(2)}$: for a weakening $\mathcal J$ of a constraint,
$\mathcal J^{(1)}$ has a solution, so by the inductive hypothesis (part (2) at
$D^{(2)}$) $\mathcal J^{(2)}$ has one. Applying the inductive hypothesis
(part (1)) at $D^{(2)}$ gives the conclusion.

\emph{Case 2: otherwise.} Lemma~\ref{lem:ubiquity} gives
$E\subt[D_{z}]{T(\sigma)}D^{(1)}_{z}$ for some $z$ and $T\in\{\TPC,\TL\}$.
Property (1a) is obtained by choosing, for each $x$, a minimal
$D^{(2)}_{x}\subte{\MT T}D^{(1)}_{x}$ containing the value at $x$ of a solution
produced by weakening a fixed constraint $C$; Lemma~\ref{lem:multitype-chain}
gives $D^{(2)}\lll^{D}D^{(1)}$ and
Lemma~\ref{lem:minimal-consistent} splits into two subcases, in one of which the
inductive hypothesis applies and in the other of which
Lemma~\ref{lem:parallelogram-crucial} applies directly.

For (1b)/(1c) one forms, for each $B$ in the family
$\mathcal B=\{B:B\subt[D_{z}]{T(\sigma)}D^{(1)}_{z}\}$, the reduction
$D^{(B,\top)}$ that equals $D^{(1)}$ off $z$ and $B$ at $z$, and the maximal
$1$-consistent reduction $D^{(B,\bot)}\subseteq D^{(B,\top)}$ (possibly empty).
Lemma~\ref{lem:tree-coverings} supplies tree-coverings $\Upsilon_{B,x}$
computing $D^{(B,\bot)}_{x}$, and their conjunction $\Upsilon_{x}$ works
uniformly in $B$. Writing $\mathcal B_{0}$ for the $B$ with $D^{(B,\bot)}$
nonempty, the argument splits:
\begin{itemize}\tightlist
\item $\mathcal B_{0}=\varnothing$: the type must be $\TL$, and
  Lemma~\ref{lem:bridge-from-instance} applied to a suitably weakened
  tree-covering produces, for any two constraints sharing a variable, a
  reflexive bridge between their $\ConOne$'s. Hence $\mathcal I$ is connected,
  and (1b) or (1c) holds according as its solution set is subdirect.
\item $\mathcal B_{0}\neq\varnothing$: one constructs a family $\Omega$ of
  weakenings of $\mathcal I$ that is \emph{critically} unable to solve all of
  $\mathcal B_{0}$ (conditions 1--4 of
  Construction~\ref{constr:omega}) and argues by whether some member fails
  (1b).
\end{itemize}
\end{proof}

\begin{hazard}[What $\Omega$'s construction turns on]\label{haz:omega}
Condition 3 is a minimality condition, and its statement is where the work is.
Read as ``if we replace any instance in $\Omega$ by \emph{all} weaker instances
then 2 is not satisfied'' it inherits the defect of
Caveat~\ref{haz:weakening}: the replacement need not make progress, and the
iteration that is supposed to produce $\Omega$ has no measure. Read as in Proposition~\ref{constr:omega} --- one member replaced by the
finitely many $\mathcal J[C\Uparrow]$, and termination measured on the
constraints rather than on the solution set --- both problems go away at once,
and the derivation of condition 4 becomes exact rather than approximate,
because what condition 3 then hands over is verbatim
Definition~\ref{def:crucial}.

Two things about the construction are worth keeping in view.

\begin{enumerate}\alphlabels\tightlist
\item \emph{The counting argument is used twice, one level apart.}
  Lemma~\ref{lem:weakening-terminates} orders \emph{instances} by counting
  their constraints by $\nu$-value; Proposition~\ref{constr:omega} orders
  \emph{sets of instances} by counting their members. Both are the multiset
  extension of a well-founded order, written as a lexicographic order on a
  product of copies of $\mathbb N$ indexed by a finite set, which is the same
  argument in a form whose well-foundedness is visible without citing
  Dershowitz--Manna.
\item \emph{Condition 4's witness depends on $\mathcal J$.} The $B$ produced
  varies with the member, so condition 4 is a $\Pi\Sigma$ statement: ``for
  every $\mathcal J$ there is $B$'' may not be strengthened to a single $B$
  serving all of $\Omega$, and the consumers in \S\ref{sec:main-statements}
  pick a fresh $B$ for each $\mathcal J$ accordingly.
\item \emph{A second descending sequence sits inside the first.} Subcase~1 of
  the $\mathcal B_{0}\neq\varnothing$ branch builds
  $\mathcal B_{1}\supseteq\mathcal B_{2}\supseteq\dots$ by repeatedly applying
  the inductive hypothesis, terminating because $\mathcal B_{1}$ is finite and
  the sequence is strictly decreasing. That is a different and much easier
  well-foundedness argument than
  Proposition~\ref{constr:omega}'s, and the two are worth keeping apart:
  nested constructions of this shape are where informal proofs most often hide
  a circularity, and neither of these does.
\end{enumerate}
\end{hazard}

\begin{longproposition}[Construction of $\Omega$]\label{constr:omega}
Let $\mathcal I$ be crucial in $D^{(1)}$ with $\mathcal B_{0}\neq\varnothing$,
and for a weakening $\mathcal J$ of $\mathcal I$ put
$\Sol(\mathcal J)=\{B\in\mathcal B_{0}:\mathcal J^{(B,\bot)}\text{ has a
solution}\}$. Then there is a finite set $\Omega$ of instances with:
\begin{enumerate}\tightlist
\item every member is a weakening of $\mathcal I$;
\item $\bigcap_{\mathcal J\in\Omega}\Sol(\mathcal J)=\varnothing$;
\item for every $\mathcal J\in\Omega$, the set
  $\bigl(\Omega\setminus\{\mathcal J\}\bigr)\cup
   \{\mathcal J[C\Uparrow]:C\in\mathcal J\}$
  does \emph{not} satisfy 2;
\item for every $\mathcal J\in\Omega$ there is $B\in\mathcal B_{0}$ with
  (a) $\mathcal J$ crucial in $D^{(B,\bot)}$ and (b) $B\in\Sol(\mathcal J')$
  for every $\mathcal J'\in\Omega\setminus\{\mathcal J\}$.
\end{enumerate}
\end{longproposition}

\begin{proof}
\emph{Monotonicity.} If $\mathcal J'$ is weaker than $\mathcal J$ then every
solution of $\mathcal J^{(B,\bot)}$ restricts to one of
$\mathcal J'^{(B,\bot)}$, so $\Sol(\mathcal J)\subseteq\Sol(\mathcal J')$.

\emph{Conditions 1--3.} Start with $\Omega=\{\mathcal I\}$. Condition 1 is
clear, and condition 2 holds because $D^{(B,\bot)}\subseteq D^{(1)}$ and
$\mathcal I^{(1)}$ has no solution, so $\Sol(\mathcal I)=\varnothing$. While
condition 3 fails at some $\mathcal J\in\Omega$ --- that is, while the
displayed set still satisfies 2 --- replace $\Omega$ by that set. Condition 1
survives, since a single weakening step applied to a weakening of $\mathcal I$
is again one; condition 2 survives by the choice of the step.

\emph{Termination.} Each $\mathcal J[C\Uparrow]$ is one weakening step from
$\mathcal J$, so by Lemma~\ref{lem:weakening-terminates}(c) it is strictly
below $\mathcal J$ in the order that lemma provides on the set $W$ of
weakenings of $\mathcal I$; and $W$ is finite by
Lemma~\ref{lem:weakening-terminates}(b). Fix a linear order on $W$ extending
it, and for a finite $\Omega\subseteq W$ let $c(\Omega)\in\mathbb N^{W}$ count
how many members $\Omega$ has at each element. One replacement removes a member
and inserts members strictly below it, so the coordinate there drops by one,
the coordinates above are untouched, and only coordinates below grow. Ordering
$\mathbb N^{W}$ lexicographically from the largest element of $W$ downwards ---
again a lexicographic order on a finite product of copies of $\mathbb N$ ---
every replacement strictly decreases $c(\Omega)$, so the loop halts at an
$\Omega$ satisfying 1--3.

This is the same counting argument twice over: once inside
Lemma~\ref{lem:weakening-terminates} to order instances by their constraints,
and once here to order sets of instances by their members.

\emph{Condition 4.} Fix $\mathcal J\in\Omega$. By condition 3 the displayed set
fails 2, so some $B\in\mathcal B_{0}$ lies in all of its members' $\Sol$. Then
$B\in\Sol(\mathcal J')$ for every $\mathcal J'\in\Omega\setminus
\{\mathcal J\}$, which is 4(b); and $B\notin\Sol(\mathcal J)$, since otherwise
$B$ would lie in $\bigcap_{\mathcal K\in\Omega}\Sol(\mathcal K)$, empty by
condition 2. So $\mathcal J^{(B,\bot)}$ has no solution while
$\mathcal J[C\Uparrow]^{(B,\bot)}$ has one for every constraint $C$ of
$\mathcal J$ --- which is Definition~\ref{def:crucial}, once no constraint of
$\mathcal J$ has a dummy variable.

None does. If $C$ had one, its projection $C'$ off that variable is weaker than
$C$ with the same violation set, and $C'$ is among the constraints of
$\mathcal J[C\Uparrow]$; so every solution of $\mathcal J[C\Uparrow]$ satisfies
$C'$ and hence $C$, giving $\Sol(\mathcal J[C\Uparrow])=\Sol(\mathcal J)$ and
so $B\in\Sol(\mathcal J)$. Finally $\mathcal J$ has at least one constraint,
since a weakening step replaces a constraint by a nonempty family rather than
deleting one, and $\mathcal I$ is crucial.
\end{proof}

\subsection{The three consumed theorems}\label{sec:codim-one}

\begin{theorem}[Reductions are safe]\label{thm:reductions-safe}
Let $\Theta$ be a cycle-consistent irreducible instance and
$B\subt[D_{x}]{T}D_{x}$ with $T\in\{\TBA,\TC,\TPC\}$. Then $\Theta$ has a
solution if and only if $\Theta$ has a solution with $x\in B$.
\end{theorem}

\begin{proof}
For $T=\TPC$ this is Theorem~\ref{thm:pc-safe}. For $T\in\{\TBA,\TC\}$,
Lemma~\ref{lem:find-consistent} gives a $1$-consistent reduction
$D^{(1)}\subte{T}D$ with $D^{(1)}_{x}\subseteq B$; by
Theorem~\ref{thm:main-inductive}(2), $\Theta^{(1)}$ has a solution, hence
$\Theta$ has one with $x\in B$.
\end{proof}

\begin{theorem}[PC reductions do not kill all solutions]\label{thm:pc-safe}
If $\mathcal I$ is cycle-consistent and irreducible, $B\subt[D_{y}]{\TPC(\sigma)}D_{y}$
for some $y\in\Var(\mathcal I)$, and $\mathcal I$ has a solution, then
$\mathcal I$ has a solution with $y\in B$.
\end{theorem}

\begin{lemma}[Subuniverses of these products are cosets]\label{lem:subuniverse-coset}
Let $q_{1},\dots,q_{r}$ be primes with $\mathbf Z_{q_{j}}\in\Vn$ and let
$\varnothing\neq L\le\mathbf Z_{q_{1}}\times\dots\times\mathbf Z_{q_{r}}$. Then
$L=g+M$ for some $g\in L$ and some subgroup $M$ of the product.
\end{lemma}

\begin{proof}
By Proposition~\ref{prop:p-divides}, $q_{j}\mid n-1$ for every $j$ and $w$ acts
on each factor as $x_{1}+\dots+x_{n}$; so $w$ acts on the product as the
$n$-fold sum, and $(n-1)h=0$ for every $h$ in the product. Fix $g\in L$ and put
$M=L-g$. If $\bar y_{1},\dots,\bar y_{n}\in M$ then
$w(\bar y_{1}+g,\dots,\bar y_{n}+g)=\sum_{j}\bar y_{j}+ng=
\bigl(\sum_{j}\bar y_{j}\bigr)+g$, which lies in $L$; so $M$ is closed under
the $n$-fold sum. Now $0\in M$, and for $y,z\in M$ taking
$(y,z,0,\dots,0)$ --- legitimate since $n\ge3$ --- gives $y+z\in M$. A finite
nonempty subset of a group closed under addition is a subgroup.
\end{proof}

\begin{longtheorem}[Codimension one]\label{thm:codim-one}
Suppose
\begin{enumerate}\romanlabels\tightlist
\item $\mathcal I$ is a linked cycle-consistent irreducible instance with
  $\Var(\mathcal I)=\{x_{1},\dots,x_{m}\}$;
\item $\mathbf D_{x_{i}}$ is $\TS$-free for every $i$;
\item $\mathcal I[C\Uparrow]$ has subdirect solution set for every constraint
  $C$ of $\mathcal I$;
\item $\sigma_{x_{i}}$ is the intersection of all linear congruences $\sigma$
  on $\mathbf D_{x_{i}}$ with $\cover{\sigma}=D_{x_{i}}^{2}$;
\item $L_{x_{i}}=\mathbf D_{x_{i}}/\sigma_{x_{i}}$;
\item $\phi\colon\mathbf Z_{q_{1}}\times\dots\times\mathbf Z_{q_{k}}\to
  L_{x_{1}}\times\dots\times L_{x_{m}}$ is a homomorphism, $q_{i}$ prime;
\item $\mathcal I[C\Uparrow]$ has a solution in $\phi(\bar a)$ for every
  constraint $C$ of $\mathcal I$ and every
  $\bar a\in\mathbf Z_{q_{1}}\times\dots\times\mathbf Z_{q_{k}}$ --- that is,
  $\mathcal I$ is crucial in every $\phi(\bar a)$
  (Definition~\ref{def:crucial}).
\end{enumerate}
Then
$\Delta=\{\bar a:\mathcal I\text{ has a solution in }\phi(\bar a)\}$ is empty,
full, or a coset of the kernel of a surjective homomorphism
$\lambda\colon\mathbf Z_{q_{1}}\times\dots\times\mathbf Z_{q_{k}}\to
\mathbf Z_{p}$ --- the solution set of the single equation
$\lambda(\bar a)=d$. Such a $\lambda$ vanishes on every factor with
$q_{i}\neq p$, so the equation involves only the coordinates over
$\mathbb F_{p}$.
\end{longtheorem}

\begin{proof}[Proof]
If $\Delta$ is full we are done. Otherwise pick $\bar b\notin\Delta$ and regard
$\phi(\bar b)$ as a reduction $D^{(1)}$ for $\mathcal I$; condition (vii) says
$\mathcal I$ is crucial in $D^{(1)}$.

\emph{Step 1: some $\ConOne(C,x)$ is a perfect linear congruence.} By
Lemma~\ref{lem:minimal-consistent}, either $C_{0}^{(1)}=\varnothing$ for some
constraint $C_{0}$, or $D^{(1)}$ is $1$-consistent.
\begin{itemize}\tightlist
\item If $C_{0}^{(1)}=\varnothing$: cruciality forces $\mathcal I=\{C_{0}\}$,
  say $C_{0}=R(y_{1},\dots,y_{t})$. By
  Lemma~\ref{lem:parallelogram-crucial}, $R$ has the parallelogram property and
  $\ConOne(R,1)$ is linear with $\cover{\ConOne(R,1)}=D_{y_{1}}^{2}$. By
  Lemma~\ref{lem:linear-on-top}, $\mathbf D_{y_{1}}/\delta\cong\mathbf Z_{p}$;
  with $\psi$ the quotient map, $\zeta=\{(a_{1},a_{2},b):\psi(a_{1})-\psi(a_{2})=b\}$
  witnesses perfection.
\item If $D^{(1)}$ is $1$-consistent: by Theorem~\ref{thm:main-inductive}(1)
  every constraint has the parallelogram property and (1b) or (1c) holds. In
  either case Lemma~\ref{lem:connected}(p) applies --- under (1c) to the
  linked connected subinstance, whose containing a constraint of $\mathcal I$
  is forced by condition (iii); under (1b) to $\mathcal I$ itself.
\end{itemize}

\emph{Step 2: add a $\mathbf Z_{p}$ coordinate.} Let $\zeta$ witness
perfection of $\ConOne(C,x)$. Add a variable $z$ with domain $\mathbf Z_{p}$,
replace $x$ in $C$ by a fresh $x'$, and add $\zeta(x,x',z)$, giving
$\mathcal I'$. Put
\[
  L=\{(\bar a,b):\mathcal I'\text{ has a solution with }z=b\text{ in }\phi(\bar a)\}
   \le\mathbf Z_{q_{1}}\times\dots\times\mathbf Z_{q_{k}}\times\mathbf Z_{p},
\]
and write $G=\mathbf Z_{q_{1}}\times\dots\times\mathbf Z_{q_{k}}$ and
$\pi\colon G\times\mathbf Z_{p}\to G$ for the projection. Condition (vii) gives
$\pi(L)=G$; in particular $L\neq\varnothing$. And $(\bar b,0)\notin L$, since a
solution of $\mathcal I'$ with $z=0$ in $\phi(\bar b)$ would, by the defining
property of $\zeta$, yield a solution of $\mathcal I$ in $\phi(\bar b)$, and
$\bar b\notin\Delta$.

\emph{Step 3: $L$ is a graph.} By Lemma~\ref{lem:subuniverse-coset},
$L=(\bar g,c)+M$ for some $(\bar g,c)\in L$ and some subgroup $M\le
G\times\mathbf Z_{p}$; and $\pi(M)=\pi(L)-\bar g=G$. The subgroup
$M\cap(\{0\}\times\mathbf Z_{p})$ is a subgroup of $\mathbf Z_{p}$, so it is
trivial or the whole of $\mathbf Z_{p}$. If it were the whole of
$\mathbf Z_{p}$ then $L$ would be closed under changing its last coordinate
arbitrarily, and $\pi(L)=G$ would then give $(\bar b,0)\in L$, which Step 2
excludes. So $M\cap(\{0\}\times\mathbf Z_{p})=0$, $\pi|_{M}$ is an isomorphism
$M\to G$, and $L$ is the graph of the affine map
\[
  f\colon G\to\mathbf Z_{p},\qquad
  f(\bar a)=c+\lambda(\bar a-\bar g),
\]
where $\lambda\colon G\to\mathbf Z_{p}$ is the homomorphism reading off the last
coordinate of $(\pi|_{M})^{-1}$.

\emph{Step 4: read off $\Delta$.} By the defining property of $\zeta$, an
$\bar a$ lies in $\Delta$ exactly when $\mathcal I'$ has a solution with $z=0$
in $\phi(\bar a)$, that is, exactly when $(\bar a,0)\in L$; so
$\Delta=f^{-1}(0)$. If $\lambda=0$ then $f$ is constant and $\Delta$ is empty
or full according as $c-\lambda(\bar g)\neq0$ or $=0$. Otherwise $\lambda$ is
onto $\mathbf Z_{p}$, since its image is a nontrivial subgroup of
$\mathbf Z_{p}$, and $\Delta$ is the coset of $\ker\lambda$ on which $f$
vanishes --- the solution set of one equation $\lambda(\bar a)=d$. Finally
$\lambda$ kills every factor $\mathbf Z_{q_{i}}$ with $q_{i}\neq p$: an element
of that factor has order dividing $q_{i}$, and its image has order dividing
$\gcd(q_{i},p)=1$. So the equation involves only the coordinates over
$\mathbb F_{p}$.
\end{proof}

\begin{hazard}[Why Step 3 is not one line]\label{haz:step2-linear-algebra}
An earlier draft compressed Steps 2--4 into ``$L$ has dimension $k$ and is the
solution set of one linear equation''. Two things make that false as stated.

$L$ is a nonempty subuniverse of an idempotent affine algebra, hence in general
an \emph{affine coset} and not a subgroup; Lemma~\ref{lem:subuniverse-coset} is
what supplies the translation, and it is where $n\ge3$ and
$q\mid n-1$ are spent. And a product of prime fields with \emph{different}
primes is not one vector space, so ``dimension $k$'' and ``codimension $1$''
have no meaning until the primary components are separated. What survives is
not a dimension count but the statement that $\Delta$ is a coset of the kernel
of a surjection onto $\mathbf Z_{p}$; that kernel has index $p$, which is the
only sense in which the codimension is one. The footnote in
\cite{ZhukSimplified} about ``$J$ may contain only variables on the same
field'' is the same issue surfacing in the algorithm, and \textup{(p3)} of
\S\ref{sec:algorithm} is organised prime by prime for exactly this reason.

What Step 2 still takes from elsewhere is $\pi(L)=G$, which is condition (vii)
read through the construction of $\mathcal I'$, and the property of $\zeta$
that relates solutions of $\mathcal I'$ with $z=0$ to solutions of
$\mathcal I$. Both belong to the half of \cite{ZhukSimplified} this document
has not yet read adversarially; see \S\ref{sec:status}.
\end{hazard}

\subsection{The supporting lemmas}

The lemmas invoked above, in dependency order.

\begin{lemma}[{\cite[Lem.~6.1]{ZhukJACM}}]\label{lem:expanded-consistency}
An expanded covering of a cycle-consistent irreducible instance is
cycle-consistent and irreducible.
\end{lemma}

\begin{lemma}\label{lem:minimal-consistent}
Suppose $D^{(1)}$ is a $1$-consistent reduction for $\mathcal I$, every
$D^{(1)}_{x}$ is $\TS$-free, $T\in\{\TPC,\TL,\TD\}$, $D^{(1)}\lll D$, and
$D^{(2)}_{x}\subte{\MT T}D^{(1)}_{x}$ is a minimal $\MT T$ subuniverse for
every $x$. Then either $C^{(2)}=\varnothing$ for some constraint $C$, or
$\mathcal I^{(2)}$ is $1$-consistent.
\end{lemma}

\begin{lemma}\label{lem:crucial-irreducible}
Let $R(x_{1},\dots,x_{n})$ be a \emph{rectangular} constraint of a
$1$-consistent instance $\mathcal I$, crucial in a reduction $D^{(\top)}$. Then
$\ConOne(R,i)$ is an irreducible congruence for every $i\in[n]$.
\end{lemma}

\begin{proof}
Take $i=1$; the other coordinates are the same. Suppose
$\ConOne(R,1)=\omega_{1}\cap\omega_{2}$ with
$\ConOne(R,1)\lneq\omega_{j}\le\mathbf D_{x_{1}}^{2}$ for $j=1,2$, and put
\[
  R_{j}(x_{1},\dots,x_{n})=\exists y\;
    \bigl(R(y,x_{2},\dots,x_{n})\wedge\omega_{j}(y,x_{1})\bigr).
\]
Each $R_{j}$ is a constraint on the same scope with
$R_{j}\supseteq R$, and the containment is proper because
$\omega_{j}\supsetneq\ConOne(R,1)$; so each $R_{j}$ is weaker than $R$ and is
among the constraints of $\mathcal I[R\Uparrow]$. And
$R_{1}\cap R_{2}=R$: a tuple in both has, for each $j$, a witness $y_{j}$ with
$R(y_{j},x_{2},\dots,x_{n})$ and $(y_{j},x_{1})\in\omega_{j}$, and rectangularity
of the first variable makes any two such witnesses $\ConOne(R,1)$-related to
each other, so $(y_{1},x_{1})\in\omega_{1}\cap\omega_{2}=\ConOne(R,1)$ and
$x_{1}$ may replace $y_{1}$.

Hence every solution of $\mathcal I[R\Uparrow]^{(\top)}$ satisfies both $R_{1}$
and $R_{2}$, so it satisfies $R$ and solves $\mathcal I^{(\top)}$ --- which has
no solution. That contradicts cruciality.
\end{proof}

\begin{hazard}[Rectangularity is a hypothesis here, not a
  conclusion]\label{haz:crucial-irreducible}
An earlier draft of this lemma read ``a constraint crucial in a reduction of a
cycle-consistent irreducible instance is rectangular, and its $\ConOne$'s are
irreducible''. That runs two statements together and gets the second one's
logic backwards: rectangularity is what the argument \emph{uses}, to know that
two witnesses for the same tuple are $\ConOne(R,1)$-related. Rectangularity of
crucial constraints is a separate fact, and it comes from the parallelogram
property (Lemma~\ref{lem:parallelogram-crucial}), not from cruciality alone.

The proof is also the reason Definition~\ref{def:crucial} quantifies over
$\mathcal I[R\Uparrow]$ rather than over single weaker constraints: what it
contradicts is that $R_{1}$ and $R_{2}$ can be imposed \emph{together}, and
each of them alone is a weakening that does gain a solution.
\end{hazard}

\begin{lemma}[A bridge from a relation]\label{lem:bridge-from-relation}
Let $R\sdle\mathbf A_{1}\times\dots\times\mathbf A_{n}$ have its first and last
variables rectangular, and suppose there are
$(b_{1},a_{2},\dots,a_{n})\in R$ and $(a_{1},\dots,a_{n-1},b_{n})\in R$ with
$(a_{1},a_{2},\dots,a_{n})\notin R$. Then
\[
  \delta(x_{1},x_{2},y_{1},y_{2})=\exists z_{2}\dots\exists z_{n-1}\;
    R(x_{1},z_{2},\dots,z_{n-1},y_{1})\wedge
    R(x_{2},z_{2},\dots,z_{n-1},y_{2})
\]
is a bridge from $\ConOne(R,1)$ to $\ConOne(R,n)$ with
$\bcol\delta=\proj_{1,n}(R)$.
\end{lemma}

\begin{proof}
Clauses (a) and (b) of Definition~\ref{def:bridge} are immediate from the
definition of $\ConOne$. For clause (d), rectangularity of the first and last
variables says exactly that $(x_{1},x_{2})\in\ConOne(R,1)$ if and only if
$(y_{1},y_{2})\in\ConOne(R,n)$ for tuples sharing the middle block
$z_{2},\dots,z_{n-1}$. Clause (c) is witnessed by
$(b_{1},a_{1},a_{n},b_{n})\in\delta$ --- take
$z_{i}=a_{i}$, so that the two conjuncts are the two tuples of the hypothesis
--- together with $(b_{1},a_{1})\notin\ConOne(R,1)$, which is what
$(a_{1},\dots,a_{n})\notin R$ says. Finally $\bcol\delta$ is $\proj_{1,n}(R)$,
since $(x,x,y,y)\in\delta$ asks for a single tuple of $R$ with first entry $x$
and last entry $y$.
\end{proof}

\begin{lemma}[Connected instances]\label{lem:connected}
Let $\mathcal I$ be a cycle-consistent connected instance. Then
\begin{enumerate}\tightlist
\item[\textup{(a)}] any two constraints with a common variable are adjacent;
\item[\textup{(b)}] for constraints $C_{1},C_{2}$, variables
  $x_{i}\in\Var(C_{i})$, and any path from $x_{1}$ to $x_{2}$, there is a
  bridge $\delta$ from $\ConOne(C_{1},x_{1})$ to $\ConOne(C_{2},x_{2})$ with
  $\bcol{\delta}$ containing all pairs connected by that path;
\item[\textup{(p)}] if $\mathcal I$ is linked then $\ConOne(C,x)$ is a perfect
  linear congruence for every $C\in\mathcal I$, $x\in\Var(C)$.
\end{enumerate}
\end{lemma}

\begin{hazard}\label{haz:connected-p}
Clause (p) is the bridge between the bridge machinery and the linear algebra;
it is what makes Theorem~\ref{thm:codim-one} possible. Its proof is (b)
followed by Lemma~\ref{lem:perfect-from-linked}: linkedness of the instance
makes $\bcol{\delta}$ linked as a relation, which upgrades the bridge. The two
senses of ``linked'' (Caveat~\ref{haz:linked}) meet exactly here, and the step
where one becomes the other should be isolated.
\end{hazard}

\begin{lemma}[{\cite[Lem.~5.6]{ZhukStrong}}]\label{lem:tree-coverings}
For a $1$-consistent instance and a reduction, there is a tree-covering whose
projection onto a given variable computes the maximal $1$-consistent
subreduction at that variable.
\end{lemma}

\begin{lemma}\label{lem:find-consistent}
Let $D^{(1)}$ be a $1$-consistent reduction of a cycle-consistent instance
$\mathcal I$ with $D^{(1)}\lll D$, and $B\subt[D_{x}]{T}D^{(1)}_{x}$ with
$T\in\{\TBA,\TC,\TPC\}$. Then there is a nonempty $1$-consistent reduction
$D^{(2)}\lll D^{(1)}$ with $D^{(2)}_{x}\subseteq B$. Moreover, if
$T\in\{\TBA,\TC\}$ then $D^{(2)}\subte{T}D^{(1)}$; and if $T=\TPC$ and every
$D^{(1)}_{y}$ is $\TS$-free then $D^{(2)}\subte{\TMPC}D^{(1)}$.
\end{lemma}

\begin{lemma}\label{lem:parallelogram-crucial}
Suppose $D^{(1)}$ is a $1$-consistent reduction for the one-constraint instance
$R(x_{1},\dots,x_{n})$, that $T\in\{\TL,\TPC,\TD\}$, that
$D^{(2)}\subte[D]{\MT T}D^{(1)}\lll D$, and that $R$ is crucial in $D^{(2)}$.
Then $R$ has the parallelogram property, and for every $i\in[n]$ the congruence
$\ConOne(R,i)$ is of type $T$ with
$\cover{\ConOne(R,i)}\supseteq(D^{(1)}_{x_{i}})^{2}$. If moreover $T=\TPC$ then
$n=2$.
\end{lemma}

\begin{longlemma}[Bridges from a subdirect PC/linear instance]\label{lem:bridge-from-instance}
Suppose
\begin{enumerate}\romanlabels\tightlist
\item $\mathcal I$ has subdirect solution set;
\item $D^{(1)}$ is a reduction with $D^{(1)}_{x}\lll D_{x}$ for every $x$;
\item $C\in\mathcal I$ is a constraint of type $T\in\{\TPC,\TL\}$;
\item $B\subt[D_{z}]{\mathcal T(\xi)}D^{(1)}_{z}$ for some variable $z$, with
  $\mathcal T\in\{\TBA,\TC,\TS,\TPC,\TL\}$;
\item if $T=\TPC$ then $\mathcal T\in\{\TPC,\TL\}$;
\item $\mathcal I^{(1)}$ has a solution;
\item $\mathcal I^{(1)}$ has no solution with $z\in B$;
\item $\mathcal I[C\Uparrow]$ has a solution in $D^{(1)}$ with $z\in B$.
\end{enumerate}
Then $\mathcal T=T$, and for every variable $x$ of $C$ there is a bridge
$\delta$ from $\xi$ to $\ConOne(C,x)$ with
$\bcol{\delta}\supseteq\mathcal I(z,x)$.
\end{longlemma}

\begin{hazard}[Stated only]\label{haz:bridge-from-instance}
Lemma~\ref{lem:bridge-from-instance} carries eight numbered hypotheses and is
\emph{stated without proof here}. It is the only source of bridges in
\S\ref{sec:main-statements}, and every use of connectedness routes through it;
so are Lemmas~\ref{lem:minimal-consistent}, \ref{lem:crucial-irreducible},
\ref{lem:bridge-from-relation}, \ref{lem:find-consistent},
\ref{lem:parallelogram-crucial}, Corollary~\ref{cor:same-type} and clauses
(a) and (b) of Lemma~\ref{lem:connected}. These are gaps in the logical chain
of this document, not imported results, and \S\ref{sec:status} records them as
such.

Its conclusion has two independent parts --- ``$\mathcal T=T$'' and ``there is a
bridge with $\bcol\delta\supseteq\mathcal I(z,x)$'' --- and different call
sites use different parts. The dependency order is
Theorem~\ref{thm:stable-intersection}, which it consumes, then this lemma, then
Theorem~\ref{thm:main-inductive}, which consumes it.
\end{hazard}

\begin{corollary}\label{cor:same-type}
In the situation of Theorem~\ref{thm:main-inductive}(2) with $\mathcal I'$
satisfying (1b), the type $\mathcal T$ of the reduction and the type of the
constraints must agree, which contradicts $\mathcal T\in\{\TBA,\TC\}$.
\end{corollary}
